% Part: first-order-logic
% Chapter: semantics
% Section: intro-semantics

\documentclass[../../../include/open-logic-section]{subfiles}

\begin{document}

\olfileid{fol}{syn}{its}

\olsection{مقدمه}

معنا بخشیدن به عبارت‌ها در قلمرو معناشناسی است. مفهوم محوری در
معناشناسی، مفهوم ارضا در !!a{structure} است. !!^a{structure} به
اجزای سازندهٔ زبان معنا می‌دهد: !!a{domain} مجموعه‌ای ناتهی از اشیاست.
سورها چنان تعبیر می‌شوند که دامنهٔ تغییرشان این دامنه باشد، به
!!{constant}s عناصری از دامنه نسبت داده می‌شود، به !!{function}s
توابعی از توانی متناهی از !!{domain} به همان دامنه نسبت داده می‌شود و
به !!{predicate}s روابطی روی !!{domain} نسبت داده می‌شود. !!{domain}
به‌همراه انتساب‌ها به واژگان پایه، !!a{structure} را تشکیل می‌دهد.
ممکن است !!^{variable}s در !!{formula}s ظاهر شوند و برای ارائهٔ
معناشناسی، باید !!{element}s از !!{domain} را نیز به آن‌ها نسبت دهیم
---این همان گمارش متغیر است. سرانجام، رابطهٔ ارضا این اجزا را به یکدیگر
پیوند می‌دهد. ممکن است !!^a{formula} در !!a{structure}~$\Struct{M}$
نسبت به گمارشِ !!a{variable}~$s$ ارضا شود؛ این امر را به‌صورت
$\Sat{M}{!A}[s]$ می‌نویسیم. این رابطه نیز با استقرا بر ساختار~$!A$
تعریف می‌شود و، برای مثال، از جدول‌های ارزش پیوندگرهای منطقی استفاده
می‌کند تا ارضای $(!A \land !B)$ را برحسب ارضاشدن (یا نشدنِ) $!A$ و
$!B$ تعریف کند. سپس روشن می‌شود که اگر !!{formula}~$!A$ یک
!!a{sentence} باشد، یعنی متغیر آزاد نداشته باشد، گمارش !!{variable}
بی‌تأثیر است؛ ازاین‌رو می‌توانیم صرفاً از ارضاشدن (یا نشدنِ)
!!{sentence}s در !!{structure}s سخن بگوییم.

بر پایهٔ رابطهٔ ارضای $\Sat{M}{!A}$ برای !!{sentence}s می‌توانیم
مفاهیم معناشناختی بنیادیِ اعتبار، استلزام و ارضاپذیری را تعریف کنیم.
!!^a{sentence} معتبر است، $\Entails !A$، اگر هر !!{structure} آن را
ارضا کند. مجموعه‌ای از !!{sentence}s آن را مستلزم است،
$\Gamma \Entails !A$، اگر هر !!{structure} که همهٔ !!{sentence}s
در~$\Gamma$ را ارضا می‌کند، $!A$ را نیز ارضا کند. همچنین مجموعه‌ای از
!!{sentence}s ارضاپذیر است اگر !!{structure}ای وجود داشته باشد که
همهٔ !!{sentence}s موجود در آن مجموعه را هم‌زمان ارضا کند. چون
!!{formula}s به‌صورت استقرایی تعریف می‌شوند و ارضا نیز به‌نوبهٔ خود
با استقرا بر ساختار !!{formula}s تعریف می‌شود، می‌توانیم برای اثبات
ویژگی‌های معناشناسی خود و مرتبط‌ساختن مفاهیم معناشناختی تعریف‌شده از
استقرا استفاده کنیم.

\end{document}
